% Supplementary Note S-Reproducibility: Step-by-Step Reproduction Guide
% Target: 3 pages; concrete commands; real expected outputs

\section*{Supplementary Note S-Reproducibility: Reproducing Paper-1 from Clone to Fresh-Eval}
\label{supp:reproducibility}

This note provides a complete command-level recipe for reproducing the three headline results of the paper from a clean environment. All commands were verified on an NVIDIA RTX~5070~Ti (16~GB) with PyTorch~2.10.0+cu128, CUDA~12.8, and Python~3.11.

\subsection*{S-R.1 Environment setup}

\begin{enumerate}[leftmargin=1.2em]
    \item Clone the repository at the canonical commit:
    \begin{verbatim}
    git clone https://github.com/USERNAME/REPO_NAME.git
    cd compute_vit
    git checkout 9cdbe77
    \end{verbatim}
    
    \item Create the conda environment:
    \begin{verbatim}
    conda env create -f environment.yml
    conda activate LLM
    \end{verbatim}
    
    \item Verify GPU and library versions:
    \begin{verbatim}
    python -c "import torch; print(torch.__version__); 
               print(torch.cuda.get_device_name(0))"
    # Expected: 2.10.0+cu128
    # Expected: NVIDIA GeForce RTX 5070 Ti
    \end{verbatim}
\end{enumerate}

\subsection*{S-R.2 Verification suite}

Run the post-fix unit-test suite to confirm the codebase is in the expected state:

\begin{verbatim}
python test_dual_bug_fix.py
# Expected: 7/7 tests passed

python test_groupwise_nl_wrapper.py
# Expected: 9/9 tests passed

python test_adc_perinstance_calibration.py
# Expected: 1/1 test passed
\end{verbatim}

If any test fails, verify that the repository is checked out at commit \texttt{9cdbe77} and that the conda environment matches \texttt{environment.yml}.

\subsection*{S-R.3 Canonical Ensemble HAT training}

Train a single canonical Ensemble HAT checkpoint (Tiny-ViT V4, uniform noise, $NL=1.0$):

\begin{verbatim}
python train_tinyvit_ensemble.py \
  --hat-training True \
  --noise-mode uniform \
  --nl-ltp 1.0 --nl-ltd -1.0 \
  --sigma-d2d 0.10 --sigma-c2c 0.05 \
  --epochs 100 --batch-size 64 --seed 123
\end{verbatim}

\textbf{Expected wall-clock:} $\sim$10--12 hours on RTX~5070~Ti.  
\textbf{Expected best test accuracy:} $\sim$91\% (single seed, training-set accuracy).  
\textbf{Output:} \texttt{checkpoints/V4\_hybrid\_standard\_noise\_hat\_best.pt}.

\textit{Note:} The training script uses automatic mixed precision (AMP) by default. Set \texttt{--amp False} only for debugging; AMP does not materially affect the final accuracy for this model size.

\subsection*{S-R.4 Fresh-instance evaluation}

Evaluate the trained checkpoint on 10 fresh D2D instances with 5 Monte-Carlo runs per instance:

\begin{verbatim}
python eval_fresh_instances_postfix.py \
  --checkpoint checkpoints/V4_hybrid_standard_noise_hat_best.pt \
  --exp-id V4 \
  --num-instances 10 \
  --mc-runs 5 \
  --nl-ltp 1.0 --nl-ltd -1.0
\end{verbatim}

\textbf{Expected wall-clock:} $\sim$30--45 minutes.  
\textbf{Expected aggregate:} $86.37 \pm 1.54\%$ (canonical Ensemble HAT fresh-instance mean across 10 instances).

The script automatically saves a JSON file with per-instance means and aggregate statistics. The output filename is derived from the checkpoint stem; verify that \texttt{fresh\_aggregate.mean} is within the $84$--$88\%$ band for a successful reproduction.

\subsection*{S-R.5 ADC ablation (hook diagnostic)}

Run the 8-bit ADC-on evaluation with per-instance calibration:

\begin{verbatim}
python scripts/_gpt/eval_fresh_instances_adc_ablation.py \
  --checkpoint checkpoints/V4_hybrid_standard_noise_hat_best.pt \
  --bits 8 \
  --calibration-scope per_instance \
  --num-instances 10
\end{verbatim}

\textbf{Expected wall-clock:} $\sim$30 minutes.  
\textbf{Expected:} ADC-on fresh-instance mean within $0.5$~pp of the ADC-off result from S-R.4, consistent with the iso-accuracy analysis in Section~\ref{subsec:iso-accuracy}.

\subsection*{S-R.6 Non-determinism note}

While all random seeds are fixed in the scripts above, some PyTorch operations (notably \texttt{cudnn.benchmark} and certain convolution backends) introduce small hardware-dependent non-determinism. The expected accuracy ranges quoted above account for this: reproduced values within $\pm 0.5$~pp of the reported means are considered successful.

To maximize determinism, set:
\begin{verbatim}
export CUBLAS_WORKSPACE_CONFIG=:4096:8
python -c "import torch; torch.use_deterministic_algorithms(True)"
\end{verbatim}
before running training or evaluation. This may reduce throughput by 5--10\%.
